% SPDX-License-Identifier: Apache-2.0
% covers: Timer
\datasheet{Timer}{Vreteno's timer: a count, a compare and an interrupt line, behind an AXI peripheral end.}

\dsfacts{\code{cpu/vreteno/src/timer.rs}}{\code{vreteno32::timer::Timer}}%
{\code{//docs:vreteno}, \code{//docs:axi}}

\subsection*{Function}

\code{Timer} is the machine timer of the Vreteno machine. It holds a
64-bit count, \code{mtime}, and a 64-bit compare, \code{mtimecmp}, and
shows each as two words on the bus. It raises \code{tirq} while the
count has reached the compare. The core reads that line as bit 7 of
\code{mip}.

It is the peripheral client of an \code{AxiPer}: requests come on
\code{req} and write beats on \code{wd}, and answers go out on
\code{ans} and \code{rb}. The input \code{rst} clears the count.

\subsection*{Parameters}

\begin{center}\footnotesize
\begin{tabular}{@{}lp{0.7\textwidth}@{}}
\toprule
Parameter & Meaning \\
\midrule
\code{I} & The width of the AXI identifier on \code{req}, \code{ans}
and \code{rb}, in bits. The tables use 2, as the board does. \\
\bottomrule
\end{tabular}
\end{center}

\subsection*{Register map}

Offsets from the timer's base. On the board the base is
\code{0x2000}, which is also \code{isa::TIMER\_BASE}.

\begin{center}\footnotesize
\begin{tabular}{@{}llp{0.55\textwidth}@{}}
\toprule
Offset & Word & Access \\
\midrule
\code{0x0} & \code{mtime}, bits 31 to 0 & Read and write. \\
\code{0x4} & \code{mtime}, bits 63 to 32 & Read and write. \\
\code{0x8} & \code{mtimecmp}, bits 31 to 0 & Read and write. \\
\code{0xc} & \code{mtimecmp}, bits 63 to 32 & Read and write. \\
\bottomrule
\end{tabular}
\end{center}

The timer selects the word by bits 3 and 2 of the address and looks at
no other bits.

\dstables{Timer}

\subsection*{Behaviour}

\begin{itemize}
\item The count adds one every cycle. With \code{rst} high it goes to
zero.
\item A read is taken when \code{rb} has room and no write is waiting
for its beat. It is answered on \code{rb} in the cycle it is taken,
with the word addressed, \code{OKAY} and \code{last} set.
\item A write is taken when no write is waiting. The word it names
goes into \code{psel} and its identifier into \code{pid}.
\item The held write's beat is taken when it has arrived and
\code{ans} has room. The lanes the strobe covers replace those bytes of
the word, the other bytes stay, and \code{ans} sends \code{OKAY}.
\item Every cycle \code{pending} takes whether \code{mtime} is at or
above \code{mtimecmp}, compared as unsigned numbers. \code{tirq} is
\code{pending}, a register's output, so the core sees in one cycle
what the timer decided in the cycle before.
\item \code{mtimecmp} starts at zero and \code{rst} does not clear it,
so the line is high until software writes a compare above the count.
\end{itemize}

\subsection*{Verification}

\begin{itemize}
\item \code{//cpu/vreteno:lockstep} hands \code{pending} to the model
for each instruction and checks \code{mtimecmp} against the model at
the halt. \code{demo\_program} sets the compare to 150 and counts the
timer's interrupt. \code{random\_programs} writes the count and the
compare with arbitrary words and runs with the timer firing.
\item The build simulates \code{Timer<2>}'s netlist against the trace
of \code{//cpu/vreteno:vreteno} under nvc and Verilator:
\code{//docs:vreteno\_sim\_timer\_timer\_tb\_test} and
\code{//docs:vreteno\_vsim\_timer\_test}.
\end{itemize}

\subsection*{Used by}

\code{Board}, as \code{timer} behind the tracker \code{ptimer}, with
its line wired to the core's \code{tirq}. The demonstration
\code{//cpu/vreteno:vreteno}, the lockstep test and \code{run.rs}.

\subsection*{Limits}

The count counts clock cycles. The timer serves bursts of one beat and
checks no address beyond bits 3 and 2, since the router sends it only
its own range. One write is held at a time, and no read is taken while
it waits. A write updates 32 bits of a 64-bit register, so software
sets a whole compare in two stores.
